\section{Background and Related Work}
\label{sec:background}

\subsection{Quantal response and its empirical content}
Logit QRE~\cite{mckelvey1995} replaces exact best response with
$\sigma_i(a)\propto\exp\{\lambda_i U_i(a;\sigma_{-i})\}$, a fixed point that interpolates
between uniform play ($\lambda\to0$) and Nash ($\lambda\to\infty$). It is exactly the
argmax of expected payoff plus $\lambda^{-1}$ times Shannon
entropy~\cite{fudenberg2015}, and $\lambda$ admits a rational-inattention reading as the
inverse shadow price of information~\cite{matejka2015}. Its principal empirical threat is
Haile, Horta\c{c}su and Kosenok~\cite{haile2008}: unrestricted quantal response rationalises
any behaviour. The standard defences are the regular-QRE axioms~\cite{goeree2016} and
external identification of payoffs. A third defence is used here and is structural rather
than parametric: the reciprocity test of Section~\ref{sec:coordinates} constrains the
\emph{symmetry} of the response matrix and admits no free parameter to absorb misfit.

\subsection{Statistical mechanics of games}
In potential games logit dynamics are reversible with a Gibbs stationary
measure~\cite{blume1993,monderer1996}; Brock and Durlauf's social-interactions
model~\cite{brock2001} is the mean-field Ising analogue and supplies the multiplicity
transition. Away from potential games the correspondence is not merely approximate---it
changes character. The chain sustains current and dissipates, and the machinery of
stochastic thermodynamics applies: Schnakenberg's cycle decomposition~\cite{schnakenberg1976},
fluctuation theorems, and the thermodynamic uncertainty relation~\cite{barato2015,seifert2012}.
Quantal response statistical equilibrium~\cite{scharfenaker2017} is the closest existing
econophysics programme; it is a constrained-MaxEnt model over aggregate outcomes and, by its
authors' own account, cannot represent out-of-equilibrium dynamics---which is precisely the
regime studied here.

\subsection{Game decomposition}
Candogan et al.~\cite{candogan2011} decompose a finite game orthogonally into potential,
harmonic and nonstrategic components, giving a basis-independent harmonic fraction $\alpha$.
Sandholm's externality symmetry~\cite{sandholm2010} and the Jacobian splitting of Balduzzi
et al.~\cite{balduzzi2018} characterise potential games; the full-versus-effective
distinction matters, and all decomposition here is taken on the normalised (strategically
equivalent) game. The Balduzzi splitting deserves a sentence of its own, because it is the
closest structural analogue to $\R$ in the games literature: decomposing the game Jacobian
into a symmetric (potential) and an antisymmetric (Hamiltonian) part is the same algebraic
move as $\chi=\tfrac12(\chi+\chi^{\top})+\tfrac12(\chi-\chi^{\top})$. Its conservation law is
Hamiltonian and deterministic; it carries no entropy production, no stationary flux and no
thermodynamics of any kind, and it never asks whether that split and a dissipation measure
are the same information. It has the local half of the pair studied here and not the global
half. Legacci, Mertikopoulos and Pradelski~\cite{legacci2024} study convergence
versus recurrence under exponential weights across the same potential--harmonic axis, in
continuous-time replicator geometry; the present work is discrete-time, works with response
matrices and stationary fluxes, and is complementary rather than competing.

\paragraph{A terminological collision, flagged because it reads like a contradiction.}
Legacci et al.\ identify the harmonic games with the \emph{incompressible} ones---zero
Shahshahani divergence of a \emph{deterministic} payoff vector field, which is a conservation
statement---while the harmonic games here are the ones with the \emph{largest} Schnakenberg
entropy production on the stochastic Glauber chain. A reader who sets ``harmonic
$\Rightarrow$ incompressible $\Rightarrow$ conservative'' beside ``harmonic $\Rightarrow$
maximally dissipative'' will conclude that the two literatures contradict each other. They do
not: one statement is about the phase-space volume of a deterministic flow, the other about
the time-asymmetry of a stochastic chain, and a divergence-free vector field implies nothing
about the reversibility of a Markov generator.

\subsection{Response and dissipation: what is already established}
\label{sec:bg:priorart}
That a response function and an entropy production are not the same information is
established physics, and it is established for precisely the class of process studied here.
Baiesi, Maes and Wynants~\cite{baiesi2009} decompose the non-equilibrium response of a Markov
jump process into a time-antisymmetric term, which is entropic and carries the dissipation,
and a time-symmetric \emph{frenetic} term built from the dynamical activity; the corollary is
that response cannot be written as a functional of entropy production alone. The revision
chain used here is a continuous-time Markov jump process, so that result applies to it
directly. \textbf{The structural half of what the $(\R,\EPR)$ plane displays is therefore
theirs and not ours}, and the same message recurs through the generalised
fluctuation--response strand~\cite{seifert2012}, which makes it settled physics rather than
one paper's position.

The converse result delimits the statement and is equally established. For Langevin dynamics,
Harada and Sasa~\cite{haradasasa2005} prove an exact equality: the frequency-integrated
violation of the fluctuation--response relation \emph{equals} the rate of energy dissipation.
In that class the local and the global reading are one number, so inequivalence is
class-dependent, and a claim of it must say why Harada--Sasa does not apply. Here it does not,
for two structural reasons rather than technical ones. First, $\R$ is the reciprocity defect
of the \emph{static} equilibrium susceptibility $\chieq=(I-SB)^{-1}S$ on the tangent
space---a property of a derivative evaluated at a fixed point---and not a time-integrated
response--correlation difference. Second, $\EPR$ is Schnakenberg's cycle functional on a
discrete profile space carrying no velocity variable, so the integral on the left-hand side of
the Harada--Sasa equality has nothing to be built from.

What this literature does not do is instantiate the pair on strategic interaction, supply a
coordinate along which the breakdown can be traced, or measure how much of the association
between the two survives once that coordinate is conditioned on. Those three things are what
this paper adds, and they are empirical claims rather than structural ones.

\subsection{Fluctuation and response in economic settings}
Garnier-Brun, Bouchaud and Benzaquen~\cite{garnierbrun2023} establish a
fluctuation--response identity for the consumer-side Slutsky matrix, which is the
demand-side counterpart of the exact static identity $\chipart=\lambda C$ used here. What
that literature has not addressed is the \emph{equilibrium} response---the object obtained
after strategic feedback---or its relationship to dissipation.

\subsection{Gap analysis}
Every ingredient exists; no tool computes them together. Gambit~\cite{gambit} is
authoritative for the QRE correspondence but exposes no derivative of the equilibrium, no
decomposition, and no dynamics. Entropy-production estimators are mature but operate on
scalar series and cannot attribute dissipation to an agent. The Candogan decomposition has
no published implementation of any kind. \textbf{The shared defect is that no existing tool
can compute a local response object and a global dissipation object for the same game at a
controlled harmonic fraction}---which is exactly the comparison this paper makes.
Table~\ref{tab:capability} states the position.

\begin{table}[H]
\caption{Capability matrix. Rows are existing tools; columns are the operations required to
compare a local response object with a global dissipation object on one game.}
\label{tab:capability}
\small
\begin{tabular}{lccccc}
\toprule
\textbf{Tool} & \textbf{QRE} & \textbf{$d\sigma^*/du$} & \textbf{Decomp.\ $\alpha$} &
\textbf{Dynamics, $\EPR$} & \textbf{Agent-aware} \\
\midrule
Gambit / pygambit~\cite{gambit} & Yes & No & No & No & --- \\
nashpy & No & No & No & No & --- \\
OpenSpiel (MMD)~\cite{sokota2023} & Partial & No & No & No & --- \\
QuantEcon (\texttt{logitdyn}) & No & No & No & Partial & Yes \\
Scalar irreversibility tools~\cite{otsubo2022} & No & No & No & Yes & \textbf{No} \\
\midrule
\strataq~\cite{strataq} & Yes & Yes & Yes & Yes & Yes \\
\bottomrule
\end{tabular}
\end{table}
